\documentclass[12pt]{article}

%----------------- Packages ------------------
\usepackage[utf8]{inputenc}
\usepackage[T1]{fontenc}
\usepackage{amsmath, amssymb}
\usepackage{geometry}
\usepackage{xcolor}
\usepackage{titlesec}
\usepackage{fancyhdr}
\usepackage[breakable]{tcolorbox}
\usepackage{graphicx}
\usepackage{tikz}
\usepackage{float}
\usetikzlibrary{arrows.meta, decorations.markings}

%----------------- Page Setup -----------------
\geometry{a4paper, margin=2.5cm}
\pagestyle{fancy}
\fancyhf{}
\rhead{Final Exam — \textbf{2022}}
\lhead{Analysis V}
\cfoot{\thepage}

%----------------- Title Styling --------------
\titleformat{\section}{\normalfont\Large\bfseries}{Exercise \thesection:}{1em}{}
\titleformat{\subsection}{\normalfont\bfseries}{Answer:}{1em}{}

%----------------- Custom Boxes ----------------
\tcbuselibrary{listingsutf8}
\newtcolorbox{answerbox}{
  colback=gray!10,
  colframe=black,
  fonttitle=\bfseries,
  title=Answer Area,
  breakable,
  before skip=10pt,
  after skip=10pt
}

%----------------- Document Start --------------
\begin{document}

%----------------- Exam Info -------------------
\begin{center}
  \Large\textbf{UNIVERSITY NAME} \\[1em]
  \large\textit{Analysis V — Final Exam} \\[0.5em]
  \large\textit{Year: 2022} \\[2em]
\end{center}

\vspace{0.5cm}

%----------------- EXERCISE 1 ------------------
\section{}
Let \( E = C^1([0,1], \mathbb{R}) \), the space of \( C^1 \) functions from \([0,1]\) to \(\mathbb{R}\). For \( f \in E \), define

\[
N(f) = \|f\|_{\infty} + \|f'\|_{\infty} \quad \text{with} \quad \|f\|_{\infty} = \sup_{x \in [0,1]} |f(x)|.
\]

\begin{enumerate}
    \item Show that \( N \) is a norm on \( E \). Is it equivalent to \( \|\cdot\|_{\infty} \)?
    
    \item 
    \begin{enumerate}
        \item Verify that \( N'(f) = |f(0)| + \|f'\|_{\infty} \) is a norm on \( E \).
        \item Show that it is equivalent to \( N \).
    \end{enumerate}
\end{enumerate}

\begin{answerbox}
\begin{enumerate}
    \item \textbf{Norm properties:}
    \begin{itemize}
        \item \textbf{Positivity}: \(N(f) \geq 0\). If \(N(f)=0\), then \(\|f\|_\infty=0\) and \(\|f'\|_\infty=0\), so \(f=0\).
        \item \textbf{Homogeneity}: \(N(\lambda f) = |\lambda|(\|f\|_\infty + \|f'\|_\infty) = |\lambda|N(f)\).
        \item \textbf{Triangle inequality}: \(N(f+g) \leq \|f\|_\infty + \|g\|_\infty + \|f'\|_\infty + \|g'\|_\infty = N(f)+N(g)\).
    \end{itemize}
    Thus \(N\) is a norm.
    
    \textbf{Equivalence with \(\|\cdot\|_\infty\)}: They are not equivalent. Consider \(f_n(x) = \sin(n\pi x)\). Then:
    \[
    \|f_n\|_\infty = 1, \quad \|f_n'\|_\infty = n\pi, \quad \text{so } N(f_n) \sim n\pi \to \infty
    \]
    while \(\|f_n\|_\infty\) remains bounded. No constant \(C\) exists such that \(N(f) \leq C\|f\|_\infty\) for all \(f\).
    
    \item 
    \begin{enumerate}
        \item \textbf{Norm properties for \(N'\)}:
        \begin{itemize}
            \item \textbf{Positivity}: Clear. If \(N'(f)=0\), then \(|f(0)|=0\) and \(\|f'\|_\infty=0\) so \(f'=0\) $\implies$ \(f\) constant, and with \(f(0)=0\) $\implies$ \(f=0\).
            \item \textbf{Homogeneity}: \(N'(\lambda f) = |\lambda f(0)| + \|\lambda f'\|_\infty = |\lambda|N'(f)\).
            \item \textbf{Triangle inequality}: \(N'(f+g) \leq |f(0)+g(0)| + \|f'+g'\|_\infty \leq N'(f)+N'(g)\).
        \end{itemize}
        
        \item \textbf{Equivalence with \(N\)}: Show \(\exists c,C>0\) such that \(cN(f) \leq N'(f) \leq CN(f)\).
        \begin{itemize}
            \item \textbf{Upper bound}: \(N'(f) = |f(0)| + \|f'\|_\infty \leq \|f\|_\infty + \|f'\|_\infty = N(f)\).
            \item \textbf{Lower bound}: For any \(x\in[0,1]\), \(f(x) = f(0) + \int_0^x f'(t)dt\), so
            \[
            |f(x)| \leq |f(0)| + \int_0^1 |f'(t)|dt \leq |f(0)| + \|f'\|_\infty.
            \]
            Taking supremum: \(\|f\|_\infty \leq |f(0)| + \|f'\|_\infty = N'(f)\).
            Thus \(N'(f) \leq N(f) \leq 2N'(f)\) (since \(N(f) = \|f\|_\infty + \|f'\|_\infty \leq N'(f) + \|f'\|_\infty \leq 2N'(f)\)).
        \end{itemize}
        So \(N\) and \(N'\) are equivalent with constants \(c=1/2\), \(C=1\).
    \end{enumerate}
\end{enumerate}
\end{answerbox}

%----------------- EXERCISE 2 ------------------
\section{}
Let \( f : \mathbb{R}^n \to \mathbb{R} \) be a continuous function and \( I \) an open interval in \( \mathbb{R} \).  
Show that \( E = \{ x \in \mathbb{R}^n : f(x) \in I \} \) is an open subset of \( \mathbb{R}^n \).

\begin{answerbox}
Let \( x_0 \in E \). Then \( f(x_0) \in I \). Since \( I \) is open, there exists \( \varepsilon > 0 \) such that
\[
(f(x_0) - \varepsilon, f(x_0) + \varepsilon) \subset I.
\]
Since \( f \) is continuous at \( x_0 \), there exists \( \delta > 0 \) such that for all \( x \in \mathbb{R}^n \) with \( \|x - x_0\| < \delta \), we have
\[
|f(x) - f(x_0)| < \varepsilon.
\]
This implies \( f(x) \in (f(x_0) - \varepsilon, f(x_0) + \varepsilon) \subset I \). Therefore, for all \( x \) in the open ball \( B(x_0, \delta) \), we have \( f(x) \in I \), i.e., \( x \in E \).

Thus, every point \( x_0 \in E \) has an open neighborhood contained in \( E \), so \( E \) is open in \( \mathbb{R}^n \).
\end{answerbox}

%----------------- EXERCISE 3 ------------------
\section{}
\begin{enumerate}
    \item Calculate
    \[
    \lim_{(x,y) \to (0,0)} \frac{\sin(xy)}{x^2 + y^2}
    \]
    
    \item Let \( f : \mathbb{R}^2 \to \mathbb{R} \) be defined by
    \[
    f(x,y) = 
    \begin{cases} 
    \dfrac{x^2 y^2}{\sqrt{(x^2 + y^2)^3}} & \text{if } (x,y) \neq (0,0) \\
    0 & \text{if } (x,y) = (0,0)
    \end{cases}
    \]
    \begin{enumerate}
        \item Show that \( \nabla f(0,0) = (0,0) \).
        \item Show that the function \( \frac{\partial f}{\partial x} : \mathbb{R}^2 \to \mathbb{R} \) is not continuous at \( (0,0) \).
    \end{enumerate}
    
    \item Let \( U \) be an open subset of \( \mathbb{R}^2 \), \( (u,v) \in U \) with \( u \neq 0 \), and \( f : U \to \mathbb{R} \) a \( C^1 \) function, homogeneous of degree \( \alpha \), such that:
    \[
    f(u,v) = 0 \quad \text{and} \quad \frac{\partial f}{\partial y}(u,v) \neq 0
    \]
    Give explicitly the implicit function \( y = \psi(x) \) that \( f \) defines near the point \( u \) and which satisfies \( \psi(u) = v \).
\end{enumerate}

\begin{answerbox}
\begin{enumerate}
    \item Use polar coordinates: \( x = r\cos\theta, y = r\sin\theta \). Then
    \[
    \frac{\sin(xy)}{x^2+y^2} = \frac{\sin(r^2\cos\theta\sin\theta)}{r^2}.
    \]
    For small \( t \), \( \sin t \sim t \), so
    \[
    \frac{\sin(r^2\cos\theta\sin\theta)}{r^2} \sim \frac{r^2\cos\theta\sin\theta}{r^2} = \cos\theta\sin\theta.
    \]
    The limit as \( r \to 0 \) depends on \( \theta \): it equals \( \cos\theta\sin\theta \), which varies with \( \theta \). For example:
    \begin{itemize}
        \item Along \( x=0 \) or \( y=0 \): limit = 0.
        \item Along \( y=x \): \( \cos\theta\sin\theta = \frac{1}{2} \) when \( \theta=\pi/4 \).
    \end{itemize}
    Since the limit depends on the direction, it does not exist.
    
    \item 
    \begin{enumerate}
        \item Compute partial derivatives at \( (0,0) \):
        \[
        \frac{\partial f}{\partial x}(0,0) = \lim_{h\to 0} \frac{f(h,0)-f(0,0)}{h} = \lim_{h\to 0} \frac{0}{h} = 0.
        \]
        Similarly \( \frac{\partial f}{\partial y}(0,0)=0 \). Thus \( \nabla f(0,0) = (0,0) \).
        
        \item For \( (x,y) \neq (0,0) \), compute \( \frac{\partial f}{\partial x} \) using quotient rule. Let \( r=\sqrt{x^2+y^2} \), then
        \[
        f(x,y) = \frac{x^2 y^2}{r^3}.
        \]
        Compute derivative (or check along a path). Consider approaching \( (0,0) \) along \( y=x \):
        For \( y=x \), \( f(x,x) = \frac{x^4}{(2x^2)^{3/2}} = \frac{x^4}{2^{3/2}|x|^3} \sim \frac{|x|}{2^{3/2}} \).
        The partial derivative \( \frac{\partial f}{\partial x} \) along this path does not tend to 0 (shows discontinuity).
        More formally: compute \( \frac{\partial f}{\partial x}(x,y) \) explicitly and show its limit as \( (x,y)\to(0,0) \) does not exist or does not equal 0.
    \end{enumerate}
    
    \item Since \( f \) is homogeneous of degree \( \alpha \), Euler's theorem gives:
    \[
    x \frac{\partial f}{\partial x}(x,y) + y \frac{\partial f}{\partial y}(x,y) = \alpha f(x,y).
    \]
    At \( (u,v) \), \( f(u,v)=0 \), so
    \[
    u \frac{\partial f}{\partial x}(u,v) + v \frac{\partial f}{\partial y}(u,v) = 0.
    \]
    By the Implicit Function Theorem, since \( \frac{\partial f}{\partial y}(u,v) \neq 0 \), there exists \( \psi \) near \( u \) with \( \psi(u)=v \) and \( f(x,\psi(x))=0 \).
    
    Differentiate \( f(x,\psi(x))=0 \):
    \[
    \frac{\partial f}{\partial x}(x,\psi(x)) + \frac{\partial f}{\partial y}(x,\psi(x)) \psi'(x) = 0.
    \]
    At \( x=u \), \( \psi(u)=v \):
    \[
    \psi'(u) = -\frac{\frac{\partial f}{\partial x}(u,v)}{\frac{\partial f}{\partial y}(u,v)}.
    \]
    Using Euler's relation: \( \frac{\partial f}{\partial x}(u,v) = -\frac{v}{u} \frac{\partial f}{\partial y}(u,v) \) (since \( u\frac{\partial f}{\partial x}+v\frac{\partial f}{\partial y}=0 \)), so
    \[
    \psi'(u) = \frac{v}{u}.
    \]
    The linear approximation near \( u \) is:
    \[
    \psi(x) \approx v + \frac{v}{u}(x-u).
    \]
    For the exact implicit function, if \( f \) is homogeneous, one can show \( \psi(x) = \frac{v}{u}x \) (if \( f \) is linear in \( y \), but generally it's the tangent approximation).
\end{enumerate}
\end{answerbox}

%----------------- END ------------------
\end{document}